\section{Related Work}

\subsection{Code and Software Engineering Benchmarks}

The evaluation of code generation capabilities has evolved from simple function-level tasks to complex repository-level challenges. \citet{chen2021evaluating} introduced HumanEval, a foundational benchmark of 164 handwritten programming problems that established the standard for measuring functional correctness in generated code. This was complemented by MBPP \citep{austin2021program}, which provided approximately 1,000 crowd-sourced Python problems designed for entry-level programmers. For more challenging algorithmic tasks, APPS \citep{hendrycks2021measuring} introduced 10,000 programming problems spanning from simple to complex algorithmic challenges.

The field has since recognized the limitations of function-level evaluation. \citet{jimenez2024swebench} pioneered repository-level evaluation with SWE-bench, presenting 2,294 real GitHub issues from 12 Python repositories that require understanding entire codebases to resolve. This revealed a significant performance gap, with state-of-the-art models resolving only the simplest issues. Building on this foundation, \citet{zan2024multiswe} extended the approach to multiple programming languages with Multi-SWE-bench, covering Java, TypeScript, JavaScript, Go, Rust, C, and C++ with 1,632 expert-curated instances. \citet{Da2025AgentRLVRTS} shows that these instances can be used for RL training as well as evaluation.

% Several benchmarks have focused on specific aspects of repository-level understanding. \citet{ding2023crosscodeeval} introduced CrossCodeEval for cross-file code completion, requiring models to leverage context from multiple files within a repository. \citet{liu2023repobench} developed RepoBench with three interconnected tasks specifically designed for evaluating repository-level auto-completion systems. More recently, \citet{zhuo2024bigcodebench} presented BigCodeBench, emphasizing code generation with diverse function calls and complex instructions. The emergence of multimodal challenges is exemplified by SWE-bench Multimodal \citep{yang2024swebenchmm}, which extends evaluation to visual software domains. These benchmarks collectively demonstrate the increasing sophistication required for comprehensive evaluation of code generation systems. \citet{He2025SWEPerfCL} explores the ability of languages models in optimizing code performance.

\subsection{Software Engineering Agents}

The development of autonomous agents capable of resolving real-world software engineering tasks has seen rapid progress. \citet{yang2024sweagent} introduced SWE-agent, emphasizing the critical importance of agent-computer interfaces (ACIs) in enabling effective code manipulation, achieving 12.5\% resolution rate on SWE-bench. This work highlighted how interface design can be as important as model capabilities for agent performance. \citet{zhang2024autocoderover} developed AutoCodeRover, which combines LLMs with sophisticated AST-based code search capabilities, achieving 19\% on SWE-bench-lite while maintaining low operational costs.

% The field has explored various architectural approaches to agent design. \citet{wang2024openhands} presented OpenHands, an open platform supporting multiple agent types and coordination mechanisms, evaluated across 15 different benchmarks. \citet{huang2023agentcoder} proposed AgentCoder, employing a multi-agent framework with specialized agents for programming, test design, and test execution, demonstrating the benefits of role specialization. \citet{wang2024executable} introduced CodeAct, which unified agent action spaces using executable Python code, showing performance improvements of up to 20\% over JSON or text-based approaches. Interestingly, \citet{xia2024agentless} challenged the complexity trend with Agentless, a simple localization-repair approach.

 